\documentclass[11pt]{article}
\usepackage[margin=1in]{geometry}
\usepackage[T1]{fontenc}
\usepackage{lmodern}
\usepackage{amsmath,amssymb}
\usepackage{booktabs}
\usepackage{threeparttable}
\usepackage[section]{placeins}  % tables stay inside their own section
\usepackage{float}              % [H]: each table stays exactly where it is input
\usepackage{graphicx}
\usepackage{xcolor}
\usepackage{caption}
\usepackage{subcaption}
\usepackage[hidelinks]{hyperref}

\captionsetup[table]{skip=4pt}
\newcommand{\note}[1]{\textcolor{gray}{\textit{[#1]}}}
\newcommand{\cmark}{\checkmark}
\newcommand{\xmark}{--}

\title{An animal identification and detection system with ImageNet-Pretrained ResNet-18}
\author{Jiacheng Chen \and Justin Chan \and Fanyou Liu \and Becky Young}
\date{}

\begin{document}
\maketitle

\begin{abstract}

\end{abstract}

\section{Introduction}

Research questions:
\begin{itemize}
  \item[RQ1] How much do ImageNet pretraining, the choice of layers to fine-tune, and augmentation help the crop
             classifier, and how does it compare with a zero-shot CLIP model and a majority-class baseline?
  \item[RQ2] How well does the detector find animals, especially small ones and several animals in one image?
  \item[RQ3] Does our detector + classifier pipeline beat a single detector that predicts the 10 classes directly?
  \item[RQ4] How robust are the models to low light, blur and low contrast situations?
\end{itemize}


\section{Related Work}
Image classification and transfer learning: ResNet~\cite{he2016deep} pretrained on ImageNet~\cite{deng2009imagenet}.\\
Object detection: Faster R-CNN~\cite{ren2015faster}, FPN~\cite{lin2017feature}; COCO
benchmark~\cite{lin2014microsoft}.\\
Zero-shot recognition: CLIP~\cite{radford2021learning}.\\
Robustness to common corruptions~\cite{hendrycks2019benchmarking}.\\

\section{Dataset}
\label{sec:dataset}
We use the animal subset of COCO train2017~\cite{lin2014microsoft} published on Hugging
Face~\cite{animaldataset}.

We split the dataset by image into training, validation and test sets (70/15/15); all experiments use this one
split (Appendix Table~\ref{tab:dataset}).

Classifier crops: each ground-truth box is expanded by 10\% per side and its short side is resized to 256 px.
Training and validation crops drop boxes with a short side under 32 px: such a box would have to be enlarged more
than six times for the classifier input, and what is left is too blurry to classify reliably. The test set keeps
every box, because small animals also appear in real photos and the detector in our pipeline passes their boxes to
the classifier; we therefore also report test accuracy by COCO size bucket (Table~\ref{tab:size}). As a result, the
classifier never sees a small box ($<32^2$ px) during training, and validation and test scores are not directly
comparable.


\section{Method}
\label{sec:method}

\subsection{Crop classifier}

\subsection{Detectors}

\subsection{Detector + classifier pipeline}

\subsection{Baselines}

\section{Experimental Setup}
\label{sec:setup}

\section{Results}
\label{sec:results}

\subsection{RQ1: Pretraining, fine-tuning and augmentation}

\begin{table}[H]
\centering
\caption{Classification of ground-truth crops (RQ1), in \%. Precision, recall and F1 are macro-averaged over the 10 classes.}
\label{tab:classification}
\begin{threeparttable}
\begin{tabular}{llcccccc}
\toprule
\multicolumn{3}{c}{Classifier} & Val & \multicolumn{4}{c}{Test} \\
\cmidrule(lr){1-3}\cmidrule(lr){4-4}\cmidrule(lr){5-8}
Init. & Trained layers & Aug. & F1 & Acc. & Prec. & Rec. & F1 \\
\midrule
ImageNet & all & \cmark & \textbf{90.2} & 78.7 & \textbf{82.6} & 81.5 & \textbf{81.6} \\
ImageNet & all & \xmark & 89.3 & \textbf{79.3} & 81.0 & 81.3 & 81.0 \\
ImageNet & head only & \cmark & 85.6 & 76.6 & 79.4 & 78.9 & 79.0 \\
Random & all & \cmark & 79.9 & 72.4 & 73.6 & 74.0 & 73.7 \\
\midrule
\multicolumn{3}{l}{CLIP zero-shot (ViT-B/32)} & -- & 78.7 & 78.6 & \textbf{83.1} & 79.5 \\
\multicolumn{3}{l}{Majority class} & -- & 16.0 & \phantom{0}1.6 & 10.0 & \phantom{0}2.8 \\
\bottomrule
\end{tabular}
\par\smallskip
{\footnotesize\noindent Validation crops exclude boxes with a short side under 32~px; test crops include them, which explains most of the gap between the validation and test F1 (Table~\ref{tab:size}).\par}
\end{threeparttable}
\end{table}

% generated by scripts/make_latex_tables.py from results/results.csv; do not edit by hand
\begin{table}[htbp]
\centering
\caption{Test accuracy by object size (COCO area buckets, pixels), in \%.}
\label{tab:size}
\begin{threeparttable}
\begin{tabular}{llcccc}
\toprule
Init. & Trained layers & Aug. & Small ($<32^2$) & Medium & Large ($\geq 96^2$) \\
\midrule
ImageNet & all & \cmark & 47.9 & \textbf{78.2} & 92.5 \\
ImageNet & all & \xmark & 52.1 & 77.4 & 92.3 \\
ImageNet & head only & \cmark & \textbf{53.4} & 71.1 & 89.7 \\
Random & all & \cmark & 52.6 & 69.9 & 82.4 \\
\midrule
\multicolumn{3}{l}{CLIP zero-shot (ViT-B/32)} & 43.2 & 76.5 & \textbf{95.5} \\
\bottomrule
\end{tabular}
\begin{tablenotes}[flushleft]\footnotesize
\item[] Every small box has a short side under 32 px, so no small box was seen during training.
\end{tablenotes}
\end{threeparttable}
\end{table}


\subsection{RQ2: Detection, small objects and several animals}

% generated by scripts/make_latex_tables.py from results/results.csv; do not edit by hand
\begin{table}[htbp]
\centering
\caption{Detection on full images (RQ2, RQ3), COCO metrics in \%.}
\label{tab:detection}
\begin{threeparttable}
\small
\begin{tabular}{lcccccccc}
\toprule
 & Val & \multicolumn{7}{c}{Test} \\
\cmidrule(lr){2-2}\cmidrule(lr){3-9}
Method & AP$_{50}$ & AP$_{50}$ & AP & AP$_{75}$ & AP$_S$ & AP$_M$ & AP$_L$ & AR$_{100}$ \\
\midrule
Animal vs.\ background$^\dagger$ & 77.8 & 76.0 & 49.4 & 54.2 & 19.3 & 41.4 & 64.4 & 58.6 \\
\midrule
One-stage, 10 classes & 75.5 & \textbf{73.8} & \textbf{48.1} & \textbf{53.6} & \textbf{15.5} & \textbf{37.2} & \textbf{58.5} & \textbf{59.5} \\
Two-stage, detector + classifier & -- & 66.8 & 44.2 & 49.2 & \phantom{0}9.6 & 32.1 & 56.2 & 54.6 \\
\bottomrule
\end{tabular}
\begin{tablenotes}[flushleft]\footnotesize
\item[] $^\dagger$Class-agnostic: a box counts if it covers any animal, so this row is not comparable with the 10-class rows below it.
\end{tablenotes}
\end{threeparttable}
\end{table}

% generated by scripts/make_latex_tables.py from results/results.csv; do not edit by hand
\begin{table}[htbp]
\centering
\caption{Test images with several animals (RQ2), in \%.}
\label{tab:multi_animal}
\begin{threeparttable}
\small
\begin{tabular}{lcccc}
\toprule
 & \multicolumn{2}{c}{$\geq$2 animals (1,603 images)} & \multicolumn{2}{c}{$\geq$2 species (233 images)} \\
\cmidrule(lr){2-3}\cmidrule(lr){4-5}
Method & Recall & Count acc. & Recall & Count acc. \\
\midrule
Animal vs.\ background & 84.7 & 35.7 & 81.7 & \textbf{29.6} \\
One-stage, 10 classes & \textbf{85.0} & 32.3 & \textbf{84.3} & 25.8 \\
Two-stage, detector + classifier & 84.8 & \textbf{38.6} & 81.8 & 29.2 \\
\bottomrule
\end{tabular}
\begin{tablenotes}[flushleft]\footnotesize
\item[] Recall: share of animals matched by a box at IoU $\geq 0.5$, whatever its label. Count accuracy: share of images where the number of boxes scoring above 0.5 equals the number of animals.
\end{tablenotes}
\end{threeparttable}
\end{table}


\subsection{RQ3: Detector + classifier vs.\ a single 10-class detector}

\begin{table}[H]
\centering
\caption{Per-class test AP$_{50}$ of the 10-class detector and the detector + classifier pipeline (RQ3), in \%.}
\label{tab:two_stage}
\begin{threeparttable}
\begin{tabular}{lccc}
\toprule
Class & 10-class detector & Detector + classifier & $\Delta$ (points) \\
\midrule
bird & \textbf{46.9} & 36.9 & $-10.0$ \\
cat & \textbf{86.1} & 80.0 & \phantom{0}$-6.1$ \\
dog & \textbf{76.9} & 67.5 & \phantom{0}$-9.4$ \\
horse & \textbf{73.9} & 68.2 & \phantom{0}$-5.7$ \\
sheep & \textbf{58.1} & 46.8 & $-11.4$ \\
cow & \textbf{59.2} & 45.2 & $-14.0$ \\
elephant & \textbf{83.9} & 80.6 & \phantom{0}$-3.3$ \\
bear & \textbf{79.5} & 75.6 & \phantom{0}$-3.9$ \\
zebra & \textbf{85.7} & 81.1 & \phantom{0}$-4.6$ \\
giraffe & \textbf{88.2} & 86.0 & \phantom{0}$-2.2$ \\
\midrule
Mean (mAP$_{50}$) & \textbf{73.8} & 66.8 & \phantom{0}$-7.1$ \\
\bottomrule
\end{tabular}
\end{threeparttable}
\end{table}


\subsection{RQ4: Robustness}

\begin{table}[H]
\centering
\caption{Robustness to image corruptions at two severities (RQ4), test set, in \%.}
\label{tab:robustness}
\begin{threeparttable}
\small
\begin{tabular}{lcccccccc}
\toprule
 & & & \multicolumn{2}{c}{Low light} & \multicolumn{2}{c}{Blur} & \multicolumn{2}{c}{Low contrast} \\
\cmidrule(lr){4-5}\cmidrule(lr){6-7}\cmidrule(lr){8-9}
Model & Metric & Clean & 1 & 2 & 1 & 2 & 1 & 2 \\
\midrule
ImageNet, all, aug. & F1 & \textbf{81.6} & \textbf{75.8} & 61.9 & \textbf{80.9} & 77.1 & \textbf{80.6} & \textbf{78.7} \\
ImageNet, all, no aug. & F1 & 81.0 & 72.6 & 60.2 & 79.7 & 75.3 & 78.7 & 75.8 \\
CLIP zero-shot & F1 & 79.5 & 72.0 & \textbf{64.4} & 79.6 & \textbf{77.7} & 77.1 & 72.5 \\
\midrule
Detector + classifier & mAP$_{50}$ & 66.8 & 60.5 & 46.8 & 59.8 & 44.8 & 64.8 & 61.7 \\
\bottomrule
\end{tabular}
\par\smallskip
{\footnotesize\noindent Low light: gamma 2 / 3 plus Gaussian noise ($\sigma$ = 0.02 / 0.04). Blur: Gaussian, radius 1 / 2~px. Low contrast: contrast factor 0.5 / 0.3. Bold marks the best classifier per column.\par}
\end{threeparttable}
\end{table}


\section{Error Analysis}

\section{Discussion and Limitations}

\section{Conclusion}

\section*{Team Contributions}

\FloatBarrier
\bibliographystyle{plain}
\bibliography{references}

\clearpage
\appendix
\section{Additional Tables}
\begin{table}[H]
\centering
\caption{Annotated boxes per class and split, and images per split. The split is by image (70/15/15, seed 42). Counts are before the 32-px filter applied to classifier training and validation crops.}
\label{tab:dataset}
\footnotesize
\setlength{\tabcolsep}{4pt}
\begin{tabular}{lcccccccccccc}
\toprule
 & \multicolumn{10}{c}{Boxes per class} & \multicolumn{2}{c}{Total} \\
\cmidrule(lr){2-11}\cmidrule(lr){12-13}
Split & bird & cat & dog & horse & sheep & cow & elephant & bear & zebra & giraffe & Boxes & Images \\
\midrule
Train & \phantom{0}7\,575 & 3\,342 & 3\,915 & 4\,740 & 6\,634 & 5\,770 & 3\,892 & \phantom{0\,}904 & 3\,712 & 3\,669 & 44\,153 & 16\,792 \\
Val & \phantom{0}1\,650 & \phantom{0\,}725 & \phantom{0\,}798 & \phantom{0\,}936 & 1\,390 & 1\,135 & \phantom{0\,}803 & \phantom{0\,}194 & \phantom{0\,}780 & \phantom{0\,}740 & \phantom{0}9\,151 & \phantom{0}3\,598 \\
Test & \phantom{0}1\,576 & \phantom{0\,}701 & \phantom{0\,}795 & \phantom{0\,}911 & 1\,485 & 1\,242 & \phantom{0\,}818 & \phantom{0\,}196 & \phantom{0\,}811 & \phantom{0\,}722 & \phantom{0}9\,257 & \phantom{0}3\,598 \\
\midrule
Total & 10\,801 & 4\,768 & 5\,508 & 6\,587 & 9\,509 & 8\,147 & 5\,513 & 1\,294 & 5\,303 & 5\,131 & 62\,561 & 23\,988 \\
\bottomrule
\end{tabular}
\end{table}

% generated by scripts/make_latex_tables.py from results/results.csv; do not edit by hand
\begin{table}[htbp]
\centering
\caption{Training setup. Both models use mixed precision on one RTX 4060 Laptop GPU (8 GB).}
\label{tab:setup}
\small
\begin{subtable}[t]{0.49\linewidth}
\centering
\caption{Crop classifiers}
\begin{threeparttable}
\begin{tabular}{@{}ll@{}}
\toprule
Model & ResNet-18 \\
Input & $224 \times 224$ crop \\
Optimizer & AdamW \\
Learning rate & $10^{-4}$ backbone, $10^{-3}$ head$^*$ \\
Weight decay & $10^{-4}$ \\
Warm-up & 1 epoch \\
Decay & cosine \\
Epochs & 15 (30 from scratch) \\
Batch size & 128 \\
Label smoothing & 0.1 \\
Sampling & class-balanced \\
Model selection & best val macro-F1 \\
\bottomrule
\end{tabular}
\begin{tablenotes}[flushleft]\footnotesize
\item[] $^*$From scratch: $10^{-3}$ for all layers.
\end{tablenotes}
\end{threeparttable}
\end{subtable}
\hfill
\begin{subtable}[t]{0.49\linewidth}
\centering
\caption{Detectors}
\begin{threeparttable}
\begin{tabular}{@{}ll@{}}
\toprule
Model & Faster R-CNN, ResNet-18 + FPN \\
Input & short side 640 px \\
Optimizer & SGD, momentum 0.9 \\
Learning rate & 0.005 \\
Weight decay & $10^{-4}$ \\
Warm-up & 500 iterations \\
Decay & $\times$0.1 at epochs 9, 12 \\
Epochs & 12 \\
Batch size & 4 \\
Frozen layers & conv1, layer1 \\
Augmentation & horizontal flip \\
Model selection & best val AP$_{50}$ \\
\bottomrule
\end{tabular}
\end{threeparttable}
\end{subtable}
\end{table}

\begin{table}[H]
\centering
\caption{Per-class test F1 of each classifier, in \%. Bold marks the best model for each class.}
\label{tab:perclass_f1}
\begin{threeparttable}
\begin{tabular}{lccccc}
\toprule
 & \multicolumn{3}{c}{ImageNet init} & Random init & CLIP \\
\cmidrule(lr){2-4}\cmidrule(lr){5-5}\cmidrule(lr){6-6}
Class & all, aug. & all, no aug. & head, aug. & all, aug. & zero-shot \\
\midrule
bird & 72.7 & \textbf{78.4} & 74.6 & 69.2 & 75.4 \\
cat & 89.8 & 89.0 & 86.4 & 76.6 & \textbf{91.6} \\
dog & \textbf{76.0} & 70.7 & 73.4 & 59.6 & 66.3 \\
horse & 78.4 & 76.4 & 74.4 & 72.3 & \textbf{78.8} \\
sheep & 71.5 & 72.4 & 67.4 & 65.8 & \textbf{74.2} \\
cow & 63.4 & 66.1 & 61.7 & 58.6 & \textbf{68.7} \\
elephant & \textbf{89.2} & 86.2 & 84.3 & 81.8 & 88.7 \\
bear & \textbf{82.8} & 78.8 & 77.4 & 68.4 & 68.3 \\
zebra & 96.1 & \textbf{96.2} & 94.7 & 91.2 & 94.7 \\
giraffe & \textbf{96.7} & 95.5 & 95.4 & 93.6 & 88.8 \\
\bottomrule
\end{tabular}
\end{threeparttable}
\end{table}


\end{document}